% Sidebar Résumé: a one-page A4 résumé with a tinted sidebar and a timeline.
% Body in Lato, name and headings in Raleway. Build with pdflatex (or latexmk).
% Edit the personal details just below, then the sidebar and main column.
\RequirePackage{fix-cm}% lets math fonts scale to any size without warnings
\documentclass[10pt]{article}% each block sets its own size below
\usepackage[a4paper,left=0pt,right=0pt,top=14mm,bottom=12mm]{geometry}
\usepackage[T1]{fontenc}
% Lato for running text: tabular lining figures keep dates and numbers aligned.
\usepackage[default,tabular,lining]{lato}
\usepackage[type1]{raleway}
\renewcommand*{\familydefault}{\sfdefault}
\usepackage[letterspace=150]{microtype}
\usepackage{xcolor}
\usepackage{tikz}
\usepackage{eso-pic}
\usepackage{fontawesome5}
\usepackage{enumitem}
\usepackage{xparse}
\usepackage[hidelinks]{hyperref}

% ---------------------------------------------------------------------------
% Personal details: the header, monogram and PDF metadata use these.
\newcommand{\cvfirst}{Elena}
\newcommand{\cvlast}{Marsh}
\newcommand{\cvinitials}{EM}
\newcommand{\cvrole}{Research Scientist \textperiodcentered{} Machine Learning for Weather \& Energy}
\hypersetup{pdftitle={\cvfirst\ \cvlast\ -- Résumé},pdfauthor={\cvfirst\ \cvlast}}

% ---------------------------------------------------------------------------
% Colours: one slate blue-grey family, plus a single warm coral detail
% (the current-position dot and the short dash under the name).
\definecolor{slate}{RGB}{62,80,102}
\definecolor{slatesoft}{RGB}{122,138,158}
\definecolor{tint}{RGB}{236,240,244}
\definecolor{hair}{RGB}{196,205,216}
\definecolor{ink}{RGB}{38,44,52}
\definecolor{coral}{RGB}{233,111,92}
% ---------------------------------------------------------------------------
% Geometry: the sidebar takes 32% of the page; both columns share one top line.
\newlength{\sidew}\setlength{\sidew}{0.32\paperwidth}
\newlength{\sidepad}\setlength{\sidepad}{8mm}
\newlength{\maingap}\setlength{\maingap}{9mm}
\newlength{\mainright}\setlength{\mainright}{13mm}

% The tinted sidebar is painted behind the page, so it always runs full height.
\AddToShipoutPictureBG{%
  \AtPageLowerLeft{\color{tint}\rule{\sidew}{\paperheight}}}

\setlength{\parindent}{0pt}
\setlength{\parskip}{0pt}
\linespread{1.12}
\frenchspacing
\pagestyle{empty}
\setlist{nosep}

% ---------------------------------------------------------------------------
% Type roles. \ralewaylt is the light display weight for the name.
\newcommand{\ralewaylt}{\fontfamily{Raleway-LF}\fontseries{l}\selectfont}
\newcommand{\ralewaysb}{\fontfamily{Raleway-LF}\fontseries{sb}\selectfont}
\newcommand{\capsl}[1]{\textls[160]{\MakeUppercase{#1}}}
\newcommand{\datefont}{\fontsize{7.6}{10}\selectfont\color{slatesoft}}
\newcommand{\metafont}{\fontsize{8.4}{11}\selectfont\color{slatesoft}}

% Section heading: letterspaced Raleway capitals with a hairline that runs out.
\newcommand{\mainsection}[1]{%
  \par\vspace{10pt}%
  {\fontsize{8.6}{10}\ralewaysb\color{slate}\capsl{#1}}%
  \hspace{2.5mm}{\color{hair}\leaders\hrule height 0.4pt depth 0pt\hfill}\kern0pt
  \par\vspace{6pt}}
\newcommand{\sidesection}[1]{%
  \par\vspace{12pt}%
  {\fontsize{8.2}{10}\ralewaysb\color{slate}\capsl{#1}}\par
  \nointerlineskip\vspace{3pt}{\color{hair}\rule{\linewidth}{0.4pt}}\par\vspace{4pt}}

% ---------------------------------------------------------------------------
% Timeline: dates on the left, a hairline with dots, the entry on the right.
% \tlentry[*]{dates}{title}{organisation, place}{details}
% The starred form ends the line at that entry (use it for the last one).
% Put \tlnow before an entry to give its dot the coral accent.
\newlength{\tldatew}\setlength{\tldatew}{17mm}
\newlength{\tlgap}\setlength{\tlgap}{6pt}
\newlength{\tlbodyw}
\newsavebox{\tlbox}
\newlength{\tldrop}
\newif\iftlnow
\newcommand{\tlnow}{\global\tlnowtrue}
\NewDocumentCommand{\tlentry}{s m m m m}{%
  \setlength{\tlbodyw}{\dimexpr\linewidth-\tldatew-7mm\relax}%
  \sbox{\tlbox}{\parbox[t]{\tlbodyw}{%
      {\bfseries #3}\par
      {\metafont #4}\par
      \ifx\relax#5\relax\else\vspace{1.5pt}{\fontsize{8.8}{11.6}\selectfont #5}\fi}}%
  \setlength{\tldrop}{\dimexpr\dp\tlbox+\tlgap+0.9ex\relax}%
  \par\noindent
  \parbox[t]{\tldatew}{\raggedleft\leavevmode\datefont #2}%
  \hspace{3.5mm}%
  \begin{tikzpicture}[overlay,baseline=0pt]
    % INVARIANT: the box is top-aligned, so its depth is everything below line one.
    \IfBooleanF{#1}{\draw[hair,line width=0.5pt]
      (0,0.9ex) -- (0,-\tldrop);}
    \iftlnow
      \fill[coral] (0,0.55ex) circle (2pt);
    \else
      \fill[white] (0,0.55ex) circle (2pt);
      \draw[slate,line width=0.7pt] (0,0.55ex) circle (1.65pt);
    \fi
  \end{tikzpicture}%
  \global\tlnowfalse
  \hspace{3.5mm}\usebox{\tlbox}\par
  \IfBooleanF{#1}{\vspace{\tlgap}}}

% ---------------------------------------------------------------------------
% Sidebar pieces.
% \contact{icon}{text}: a fixed icon column so the text lines up.
\newcommand{\contact}[2]{%
  \par\noindent\makebox[5mm][l]{\color{slate}\footnotesize #1}%
  \parbox[t]{\dimexpr\linewidth-5mm}{\raggedright #2}\par\vspace{3pt}}
% \skill{name}{level 0..1}: a thin proficiency bar under the name.
\newcommand{\skill}[2]{%
  \par\noindent{\fontsize{8.6}{10}\selectfont #1}\par
  % The bar hangs from the label, not from the baseline grid.
  \nointerlineskip\vspace{3.5pt}%
  \noindent\begin{tikzpicture}
    \fill[hair] (0,0) rectangle (\linewidth,1.1pt);
    \fill[slate] (0,0) rectangle ({#2*\linewidth},1.1pt);
  \end{tikzpicture}\par\vspace{6.5pt}}
% \lang{name}{level}: name left, level right in quiet small text.
\newcommand{\lang}[2]{%
  \par\noindent{\fontsize{8.6}{10}\selectfont #1}\hfill{\metafont #2}\par\vspace{3pt}}

% \award{year}{name}: tabular year in a narrow column, then the award.
\newcommand{\award}[2]{%
  \par\noindent\makebox[9mm][l]{\datefont #1}%
  \parbox[t]{\dimexpr\linewidth-9mm}{\raggedright #2}\par\vspace{3pt}}

% Monogram placeholder: replace with \includegraphics if you use a photo.
\newcommand{\monogram}{%
  \begin{tikzpicture}
    \fill[white] (0,0) circle (13.5mm);
    \draw[slatesoft,line width=0.4pt] (0,0) circle (13.5mm);
    \fill[slate] (0,0) circle (12mm);
    \node[white,font=\fontsize{24}{24}\ralewaylt] at (0,0)
      {\textls[60]{\cvinitials}};
  \end{tikzpicture}}

\begin{document}
\noindent\hspace{\sidepad}%
% ============================== SIDEBAR ===================================
\begin{minipage}[t]{\dimexpr\sidew-2\sidepad}
  \fontsize{8.8}{12}\selectfont\raggedright\color{ink}
  \vspace{0pt}
  \centering\monogram\par
  \raggedright

  \sidesection{Contact}
  \contact{\faEnvelope[regular]}{\href{mailto:elena.marsh@example.org}{elena.marsh@example.org}}
  \contact{\faPhone}{+44 7700 900 418}
  \contact{\faMapMarker*}{Bristol, United Kingdom}
  \contact{\faGlobe}{\href{https://elenamarsh.example.org}{elenamarsh.example.org}}

  \sidesection{Skills}
  \skill{Probabilistic forecasting}{0.95}
  \skill{Python, PyTorch, JAX}{0.92}
  \skill{Spatio-temporal deep learning}{0.88}
  \skill{Data pipelines (Dask, Zarr, Airflow)}{0.80}
  \skill{Cloud and MLOps (AWS, Docker)}{0.68}
  \skill{Technical writing and teaching}{0.90}

  \sidesection{Languages}
  \lang{English}{native}
  \lang{German}{fluent \textperiodcentered{} C1}
  \lang{Spanish}{working \textperiodcentered{} B2}

  \sidesection{Links}
  \contact{\faGithub}{\href{https://github.com/emarsh}{github.com/emarsh}}
  \contact{\faOrcid}{\href{https://orcid.org/0000-0002-1825-0097}{0000-0002-1825-0097}}
  \contact{\faLinkedinIn}{\href{https://linkedin.com/in/elenamarsh}{in/elenamarsh}}
  \contact{\faGraduationCap}{Google Scholar: E.~Marsh}

  \sidesection{Awards}
  \award{2024}{RMetS L.\,F. Richardson Prize}
  \award{2022}{Turing Institute Fellowship}
  \award{2020}{LMU Dissertation Prize}

  \sidesection{Interests}
  Mountain hydrology, open data for public weather services, mentoring
  first-generation students, long-distance cycling.
\end{minipage}%
\hspace{\dimexpr\sidepad+\maingap}%
% ============================== MAIN COLUMN ===============================
\begin{minipage}[t]{\dimexpr\paperwidth-\sidew-\maingap-\mainright}
  \fontsize{9.3}{12.6}\selectfont\color{ink}
  \vspace{0pt}
  % Name: large light Raleway, the last name a step darker for rhythm.
  {\fontsize{30}{32}\ralewaylt\color{slate}\textls[30]{\cvfirst}~{\ralewaysb\textls[30]{\cvlast}}}\par
  \vspace{5pt}
  {\color{coral}\rule{9mm}{1.4pt}}\par
  \vspace{5pt}
  {\fontsize{10.4}{13}\ralewaylt\color{ink}\cvrole}\par
  \vspace{8pt}

  \mainsection{Profile}
  Atmospheric scientist turned machine-learning researcher. I build
  probabilistic forecasting models that turn reanalysis and satellite data
  into decisions for wind and solar operators, and I care about models
  that are calibrated, explainable and cheap enough to run every hour.
  Research posts since 2020 in a university lab and an energy company; looking
  for a role that keeps one foot in open research.

  \mainsection{Experience}
  \tlnow
  \tlentry{2023 -- now}{Senior Research Scientist}{Northwind Forecasting Ltd, Bristol}{%
    Lead a team of four building the company's day-ahead wind power
    ensemble; cut the 24-hour error by 18\% against the previous system
    and brought inference on 1\,200 farms under six minutes.}
  \tlentry{2020 -- 2023}{Postdoctoral Researcher}{School of Geographical Sciences, University of Bristol}{%
    Developed graph neural networks for sub-seasonal precipitation;
    co-supervised three MSc theses and taught the module \emph{Data
    Science for the Environment}.}
  \tlentry*{2019}{Research Intern}{European Centre for Medium-Range Weather Forecasts, Reading}{%
    Prototyped neural post-processing of 2\,m temperature ensembles;
    the code fed the centre's first machine-learning pilot.}

  \mainsection{Education}
  \tlentry{2015 -- 2020}{PhD in Atmospheric Physics}{Ludwig-Maximilians-Universität München}{%
    Thesis: \emph{Learning the error of convection-permitting forecasts}.
    Summa cum laude; DWD doctoral fellowship.}
  \tlentry*{2013 -- 2015}{MSc in Physics of the Earth}{University of Edinburgh}{}

  \mainsection{Selected projects}
  \textbf{skillcurve} \quad{\metafont open-source, 2.1k GitHub stars}\par
  A Python library for verifying probabilistic forecasts (CRPS, reliability,
  sharpness) on labelled arrays; used by two national weather services.\par
  \vspace{5pt}
  \textbf{Open Solar Nowcast} \quad{\metafont Turing Institute data challenge, 2022}\par
  Organised an open benchmark for 15-minute solar nowcasting from
  satellite data; 140 teams from 30 countries took part.

  \mainsection{Selected publications}
  \begin{list}{}{\leftmargin=4mm\labelwidth=3mm\labelsep=1mm\itemsep=3.5pt\parsep=0pt\topsep=0pt}
    \item[\color{slatesoft}\textbullet] \textbf{Marsh, E.}, Okoye, T.\ and Lund, P.
      Calibrated graph networks for sub-seasonal rainfall.
      \emph{Geophysical Research Letters}, 51, 2024.
    \item[\color{slatesoft}\textbullet] \textbf{Marsh, E.} and Feld, J.
      Learning forecast error at the convective scale.
      \emph{Quarterly Journal of the Royal Meteorological Society}, 147, 2021.
    \item[\color{slatesoft}\textbullet] Haddad, R., \textbf{Marsh, E.} \emph{et al.}
      Ensemble wind power forecasting at fleet scale.
      \emph{NeurIPS Workshop on Climate Change AI}, 2023.
  \end{list}
\end{minipage}
\end{document}
